\documentclass[12pt, a4paper]{report}
\usepackage[polish]{babel}
\usepackage[utf8]{inputenc}
\usepackage{graphicx,wrapfig, lipsum}
\usepackage{a4wide}
\usepackage{titlesec}
\usepackage[T1]{fontenc}
\usepackage{adjustbox}
\usepackage{multirow}
\usepackage[table,xcdraw]{xcolor}
\usepackage{hyperref}
\usepackage{enumitem}
\usepackage{float}
\usepackage{geometry}
\usepackage{xcolor}
\usepackage[cache=false]{minted}
\usepackage{listings}
\usepackage{amsmath} 
\usepackage{amssymb}
\usepackage{xurl}
\usepackage{chngcntr}

\counterwithin{figure}{section}
\makeatletter
\renewcommand{\thefigure}{\thesection\arabic{figure}}

\usemintedstyle{autumn}

\newcommand {\R}{\mathbb{R}}
\newcommand {\C}{\mathbb{C}}
\newcommand {\N}{\mathbb{N}}

\geometry{
    left = 25mm,
    right = 25mm,
    top = 25mm,
    bottom = 30mm
}

\renewcommand{\thesection}{\arabic{section}.}
\renewcommand{\thesubsection}{\arabic{section}.\arabic{subsection}.}

\titleformat*{\section}{\LARGE\bfseries}
\renewcommand{\figurename}{Rys}
\setlength{\intextsep}{2cm}
\setlength{\textfloatsep}{0cm}

\begin{document}
\title{
  \Huge \textbf{Eksploracja Danych - projekt no.1} 
  \\ \vspace{10pt} 
  \LARGE  Pokemon- Weedle's Cave
}

\author{Alicja Sobiech, 188861\\
Nikolai Lavrinov, 201302\\
Maciej Łukasiewicz, 197865}


\maketitle

% ----------------------------------- SPIS TRESCI ----------------------------------

\tableofcontents

% ---------------------------------------------------------------------------------
% PAPER ITSELF --------------------------------------------------------------------
% ---------------------------------------------------------------------------------

\newpage
\addcontentsline{toc}{part}{Sprawozdanie}

% okreslenie celu eksploracji i kryteriow sukcesu
\section{Cel eksploracji} 

Celem eksploracji jest ....

\section{Zbiór danych}

% Ogólny opis zbioru danych
Zbiór danych zawiera ....

\subsection{Charakterystyka zbioru danych}
\begin{itemize}
    \item \textbf{Pochodzenie}\\ 
    \url{https://www.kaggle.com/datasets/terminus7/pokemon-challenge}
    
    \item \textbf{Format}\\
    
    \item \textbf{Liczba przykładów}\\
    
    \item \textbf{Ilość zbiorów danych}\\ 

\end{itemize}

\subsection{Opis atrybutów}
% opis atrybutow (nazwy, typy - nominalne/numeryczne, co oznaczaja, co oznaczaja wartosci - jednostka miary, wartosci specjalne)

\renewcommand{\figurename}{Tabela.}
\begin{adjustbox}{center, caption={table description etc},nofloat=figure,vspace=\bigskipamount}
% PLACE FOR A TABLE
\end{adjustbox}

\section{Eksploracyjna Analiza Danych}

\subsection{Wstępne ustalenia dotyczące zawartości zbioru}

\subsection{Wyniki eksploracyjnej analizy danych}

\subsubsection{Wyniki analizy danych a cel eksploracji}

\subsubsection{Rozkłady wartości atrybutów}
\subsubsection{Korelacja pomiędzy wartościami atrybutów}


\subsection{Uwagi na temat jakości danych}
\subsubsection{Dane brakujące}
\subsubsection{Punkty oddalone}
\subsubsection{Dane niespójne}
\subsubsection{Dane niezrozumiałe}

\section{Podsumowanie}


% ---------------------------------------------------------------------------------
% SOURCES -------------------------------------------------------------------------
% ---------------------------------------------------------------------------------
\newpage
\addcontentsline{toc}{part}{Źródła}

\addcontentsline{toc}{section}{Wykaz Literatury}

\section*{Wykaz Literatury}
\renewcommand{\labelenumi}{[\arabic{enumi}]}
\begin{enumerate}
\item 

\end{enumerate}


\end{document}
